% META: {"id": "1SPE-DERIVATION-GLOBAL-QCM", "chapitre": "1SPE-DERIVATION-GLOBAL", "type_objet": "qcm", "genere_depuis": "chapitres/1SPE-DERIVATION-GLOBAL/qcm/1SPE-DERIVATION-GLOBAL-QCM.json", "status": "generated"}
% Fichier genere par scripts/build_qcm_tex.py — ne pas editer a la main.

\section*{\textcolor{chapcolor}{\MakeUppercase{Fonction dérivée et applications — QCM d'auto-évaluation}}}

\begin{center}
\textit{Pour chaque question, une seule reponse est exacte.}
\end{center}

\begin{enumerate}

\item[] \textbf{Capacite C1}

\item \textbf{[Q1]} La dérivée de la fonction $f(x) = x^4$ sur $\mathbb{R}$ est :
  \begin{enumerate}[label=\Alph*.]
    \item $x^3$
    \item $4x^3$
    \item $4x^4$
    \item $3x^3$
  \end{enumerate}

\item \textbf{[Q2]} La fonction $f(x) = \sqrt{x}$ est dérivable sur $]0 ; +\infty[$ et sa dérivée est :
  \begin{enumerate}[label=\Alph*.]
    \item $\dfrac{2}{\sqrt{x}}$
    \item $\dfrac{1}{\sqrt{x}}$
    \item $\dfrac{1}{2\sqrt{x}}$
    \item $\dfrac{\sqrt{x}}{2}$
  \end{enumerate}

\item \textbf{[Q3]} Pour $x \neq 0$, la dérivée de $f(x) = \dfrac{1}{x}$ est :
  \begin{enumerate}[label=\Alph*.]
    \item $\dfrac{1}{x^2}$
    \item $-x^2$
    \item $-\dfrac{1}{x^2}$
    \item $\dfrac{1}{x}$
  \end{enumerate}

\item[] \textbf{Capacite C2}

\item \textbf{[Q4]} La dérivée de la fonction $f(x) = 3x^2 - 5x + 2$ sur $\mathbb{R}$ est :
  \begin{enumerate}[label=\Alph*.]
    \item $6x - 5$
    \item $3x - 5$
    \item $6x^2 - 5$
    \item $6x - 3$
  \end{enumerate}

\item \textbf{[Q5]} La formule de dérivation d'un produit $(uv)'$ est :
  \begin{enumerate}[label=\Alph*.]
    \item $u'v'$
    \item $u' + v'$
    \item $u'v + uv'$
    \item $u'v - uv'$
  \end{enumerate}

\item \textbf{[Q6]} La dérivée du quotient $\left(\dfrac{u}{v}\right)'$ est égale à :
  \begin{enumerate}[label=\Alph*.]
    \item $\dfrac{u'v + uv'}{v^2}$
    \item $\dfrac{u'v - uv'}{v^2}$
    \item $\dfrac{u'}{v'}$
    \item $\dfrac{u'v - uv'}{v}$
  \end{enumerate}

\item[] \textbf{Capacite C3}

\item \textbf{[Q7]} Si $f'(x) > 0$ sur un intervalle $I$, alors la fonction $f$ est :
  \begin{enumerate}[label=\Alph*.]
    \item strictement décroissante sur $I$
    \item constante sur $I$
    \item strictement croissante sur $I$
    \item nulle sur $I$
  \end{enumerate}

\item \textbf{[Q8]} Si le tableau de signe de $f'(x)$ indique $f'(x) < 0$ sur $]1 ; 4[$, alors $f$ est :
  \begin{enumerate}[label=\Alph*.]
    \item croissante sur $]1 ; 4[$
    \item nulle en $x=2{,}5$
    \item strictement décroissante sur $]1 ; 4[$
    \item convexe sur $]1 ; 4[$
  \end{enumerate}

\item \textbf{[Q9]} Si $f'(x)$ s'annule en $x=a$ en passant du signe $+$ au signe $-$, la fonction $f$ admet en $a$ :
  \begin{enumerate}[label=\Alph*.]
    \item un minimum local
    \item un point d'inflexion
    \item un maximum local
    \item une asymptote verticale
  \end{enumerate}

\item[] \textbf{Capacite C4}

\item \textbf{[Q10]} Soit $B(x) = -x^2 + 6x - 5$ la fonction bénéfice. Le bénéfice maximal est atteint pour $x$ égal à :
  \begin{enumerate}[label=\Alph*.]
    \item $-3$
    \item $3$
    \item $6$
    \item $0$
  \end{enumerate}

\item \textbf{[Q11]} Pour une fonction polynôme du second degré $f(x) = ax^2 + bx + c$ avec $a < 0$, l'extremum est :
  \begin{enumerate}[label=\Alph*.]
    \item un minimum global
    \item un maximum global
    \item un point d'inflexion
    \item inexistant
  \end{enumerate}

\item \textbf{[Q12]} Pour déterminer les extremums locaux d'une fonction dérivable $f$ sur un intervalle ouvert, on cherche :
  \begin{enumerate}[label=\Alph*.]
    \item les valeurs où $f(x) = 0$
    \item les valeurs où $f'(x) = 0$ avec changement de signe
    \item les valeurs où $f''(x) > 0$
    \item les bornes de l'intervalle
  \end{enumerate}

\item[] \textbf{Capacite C5}

\item \textbf{[Q13]} Dans un problème d'optimisation, la première étape consiste à :
  \begin{enumerate}[label=\Alph*.]
    \item dériver immédiatement
    \item exprimer la quantité à optimiser sous la forme d'une fonction $f(x)$ d'une seule variable
    \item résoudre l'équation $f(x) = 0$
    \item tracer la courbe sur la calculatrice
  \end{enumerate}

\item \textbf{[Q14]} Un fabricant veut maximiser le volume $V(x) = x(120 - 2x)^2$ pour $x \in ]0 ; 60[$. $V'(x) = 0$ donne $x$ égal à :
  \begin{enumerate}[label=\Alph*.]
    \item $60$
    \item $40$
    \item $20$
    \item $10$
  \end{enumerate}

\item \textbf{[Q15]} Pour trouver le maximum d'une fonction continue $f$ sur un segment fermé $[a, b]$, il faut comparer :
  \begin{enumerate}[label=\Alph*.]
    \item les valeurs de $f'$ aux bornes
    \item les valeurs de $f$ aux points critiques intérieurs où $f'(x)=0$ et aux bornes $f(a), f(b)$
    \item uniquement les points où $f'(x)=0$
    \item uniquement les bornes $f(a)$ et $f(b)$
  \end{enumerate}

\end{enumerate}

% NEXUS-QCM-TEACHER-ONLY-BEGIN
\ifnxVersionProfesseur
\clearpage
\section*{Cle de correction — reservee au professeur}

\begin{center}
\begin{tabular}{lll}
\hline
Question & Capacite & Reponse exacte \\
\hline
Q1 & C1 & \textbf{B} \\
Q2 & C1 & \textbf{C} \\
Q3 & C1 & \textbf{C} \\
Q4 & C2 & \textbf{A} \\
Q5 & C2 & \textbf{C} \\
Q6 & C2 & \textbf{B} \\
Q7 & C3 & \textbf{C} \\
Q8 & C3 & \textbf{C} \\
Q9 & C3 & \textbf{C} \\
Q10 & C4 & \textbf{B} \\
Q11 & C4 & \textbf{B} \\
Q12 & C4 & \textbf{B} \\
Q13 & C5 & \textbf{B} \\
Q14 & C5 & \textbf{C} \\
Q15 & C5 & \textbf{B} \\
\hline
\end{tabular}
\end{center}

\subsection*{Diagnostic des reponses erronees}

\begin{itemize}
  \item \textbf{Q1}
  \begin{itemize}
    \item \textbf{A} — Oubli du coefficient $n$ dans la formule $(x^n)' = n x^{n-1}$. \emph{Renvoi : M1.}
    \item \textbf{C} — Exposant non décrémenté ($4x^4$ au lieu de $4x^3$). \emph{Renvoi : M1.}
    \item \textbf{D} — Coefficient $n$ remplacé par $n-1 = 3$. \emph{Renvoi : M1.}
  \end{itemize}
  \item \textbf{Q2}
  \begin{itemize}
    \item \textbf{A} — Facteur 2 au numérateur au lieu du dénominateur. \emph{Renvoi : M1.}
    \item \textbf{B} — Dénominateur 2 omis dans la formule de dérivation de $\sqrt{x}$. \emph{Renvoi : M1.}
    \item \textbf{D} — $\sqrt{x}$ divisé par 2 sans dérivation. \emph{Renvoi : M1.}
  \end{itemize}
  \item \textbf{Q3}
  \begin{itemize}
    \item \textbf{A} — Signe moins oublié dans la dérivée de $1/x$. \emph{Renvoi : M1.}
    \item \textbf{B} — Forme $-x^2$ au lieu de $-1/x^2$. \emph{Renvoi : M1.}
    \item \textbf{D} — L'eleve recopie la fonction $1/x$ sans la deriver ; sa derivee vaut $-1/x^2$. \emph{Renvoi : M1.}
  \end{itemize}
  \item \textbf{Q4}
  \begin{itemize}
    \item \textbf{B} — Coefficient 3 non dérivé, terme $3x$ au lieu de $6x$. \emph{Renvoi : M2.}
    \item \textbf{C} — Exposant non décrémenté. \emph{Renvoi : M2.}
    \item \textbf{D} — Terme constant $+2$ conservé dans la dérivée. \emph{Renvoi : M2.}
  \end{itemize}
  \item \textbf{Q5}
  \begin{itemize}
    \item \textbf{A} — Produit des dérivées au lieu de la formule Leibniz. \emph{Renvoi : M2.}
    \item \textbf{B} — Somme des dérivées. \emph{Renvoi : M2.}
    \item \textbf{D} — Signe moins (c'est le numérateur du quotient). \emph{Renvoi : M2.}
  \end{itemize}
  \item \textbf{Q6}
  \begin{itemize}
    \item \textbf{A} — Signe $+$ au numérateur au lieu de $-$. \emph{Renvoi : M2.}
    \item \textbf{C} — Quotient des dérivées. \emph{Renvoi : M2.}
    \item \textbf{D} — Carré au dénominateur oublié. \emph{Renvoi : M2.}
  \end{itemize}
  \item \textbf{Q7}
  \begin{itemize}
    \item \textbf{A} — Signe $+$ de $f'$ confondu avec décroissance. \emph{Renvoi : M3.}
    \item \textbf{B} — Signe $+$ de $f'$ confondu avec $f'$ nulle. \emph{Renvoi : M3.}
    \item \textbf{D} — Signe de $f'$ confondu avec annulation de $f$. \emph{Renvoi : M3.}
  \end{itemize}
  \item \textbf{Q8}
  \begin{itemize}
    \item \textbf{A} — Intervalle de croissance pris au lieu de décroissance. \emph{Renvoi : M3.}
    \item \textbf{B} — Racine de $f$ au lieu du signe de $f'$. \emph{Renvoi : M3.}
    \item \textbf{D} — Le signe de $f'$ determine les variations de $f$, pas sa convexite. La convexite s'etudie a l'aide de $f''$. \emph{Renvoi : M3.}
  \end{itemize}
  \item \textbf{Q9}
  \begin{itemize}
    \item \textbf{A} — Changement $+ \to -$ correspond à un maximum, pas un minimum. \emph{Renvoi : M3, M4.}
    \item \textbf{B} — Annulation de $f'$ avec changement de signe donne un extremum. \emph{Renvoi : M4.}
    \item \textbf{D} — Une asymptote verticale se detecte par le comportement limite de $f$ au voisinage de $a$, pas par l'annulation de $f'$ en $a$. \emph{Renvoi : M3.}
  \end{itemize}
  \item \textbf{Q10}
  \begin{itemize}
    \item \textbf{A} — Signe de $-2x$ oublié : $-2x+6=0$ donne $x=3$, pas $-3$. \emph{Renvoi : M4.}
    \item \textbf{C} — $x=6$ correspond à $B'(6) = -6 \neq 0$. \emph{Renvoi : M4.}
    \item \textbf{D} — $x=0$ correspond à $B'(0) = 6 \neq 0$. \emph{Renvoi : M4.}
  \end{itemize}
  \item \textbf{Q11}
  \begin{itemize}
    \item \textbf{A} — Coefficient dominant négatif implique un maximum global, pas un minimum. \emph{Renvoi : M4.}
    \item \textbf{C} — Pas de point d'inflexion pour une parabole. \emph{Renvoi : M4.}
    \item \textbf{D} — $f'$ change de signe donc c'est bien un extremum. \emph{Renvoi : M4.}
  \end{itemize}
  \item \textbf{Q12}
  \begin{itemize}
    \item \textbf{A} — $f(x)=0$ donne les racines de $f$, pas ses extremums. \emph{Renvoi : M4.}
    \item \textbf{C} — Condition de convexité, pas seule d'extremum. \emph{Renvoi : M4.}
    \item \textbf{D} — Sur un intervalle ouvert, les bornes ne sont pas atteintes. \emph{Renvoi : M3.}
  \end{itemize}
  \item \textbf{Q13}
  \begin{itemize}
    \item \textbf{A} — Dériver avant d'avoir défini la fonction à optimiser. \emph{Renvoi : M5.}
    \item \textbf{C} — Résoudre $f(x)=0$ cherche les racines, pas l'optimum. \emph{Renvoi : M5.}
    \item \textbf{D} — Tracer la courbe n'est pas une preuve algébrique. \emph{Renvoi : M5.}
  \end{itemize}
  \item \textbf{Q14}
  \begin{itemize}
    \item \textbf{A} — $V'(60)=0$, mais $60$ est hors de l'intervalle ouvert $]0;60[$ impose par l'enonce ; ce n'est donc pas une valeur admissible de $x$. \emph{Renvoi : M5.}
    \item \textbf{B} — L'eleve a remplace a tort le facteur $120 - 6x$ de $V'(x)=(120-2x)(120-6x)$ par $120 - 3x$ ; resoudre $120 - 3x=0$ conduit exactement a $x=40$. \emph{Renvoi : M5.}
    \item \textbf{D} — L'eleve a remplace a tort le facteur $120 - 6x$ par $60 - 6x$ ; resoudre $60 - 6x=0$ conduit exactement a $x=10$. \emph{Renvoi : M5.}
  \end{itemize}
  \item \textbf{Q15}
  \begin{itemize}
    \item \textbf{A} — Les valeurs de $f'$ ne sont pas les valeurs de la fonction $f$. \emph{Renvoi : M5.}
    \item \textbf{C} — Sur un segment fermé, le maximum peut être atteint sur une borne sans annulation de $f'$. \emph{Renvoi : M5.}
    \item \textbf{D} — Le point critique intérieur doit aussi être vérifié. \emph{Renvoi : M5.}
  \end{itemize}
\end{itemize}
\fi
% NEXUS-QCM-TEACHER-ONLY-END
